
\subsubsection{Methodological Contribution}

The central contribution is a validation framework that distinguishes statistical independence from directional factorization. Statistical independence means metrics are uncorrelated (low coupling). Directional factorization means code changes align along metric-specific axes (spectral gap > 2.0, informative labels). Our synthetic data demonstration shows these can be dissociated.

\textbf{Potential architectural implication (if validated on real data)}: Orthogonality constraints in multi-objective architectures may be redundant when data is already independent. These constraints prevent task interference by enforcing geometric separation—but if metrics are naturally uncorrelated, interference is minimal. The complexity of factorized architectures would lack empirical justification when covariance geometry is spherical.

\textbf{Methodological implication}: Researchers proposing factorized architectures should validate empirical separability first using methods like those presented here, rather than assuming structure transfers from vision/NLP multi-task learning.

\subsubsection{Why Test This?}

Several factors suggest this warrants empirical validation:

\textbf{Intent-Outcome Gap Hypothesis}: Commit messages may reflect developer intent ("fix security vulnerability"), but actual outcomes could be multi-faceted. A security fix may add input validation (affecting quality and efficiency), refactor error handling (affecting correctness and maintainability), or update dependencies (affecting multiple dimensions unpredictably).

\textbf{Measurement Tool Independence}: The independence in synthetic data (coupling = 0.072) reflects measurement design—SonarQube measures different things than CodeQL measures different things than pytest-benchmark. Tools are independent by construction. Behavioral factorization would require developers to use this independence—making changes that affect one tool dominantly.

\textbf{Code Structure Constraints}: Code modifications may be inherently coupled. Improving security often requires refactoring (quality), which may introduce complexity (efficiency trade-off), which needs testing (correctness verification).

\subsubsection{Limitations}

\textbf{L1: Synthetic Test Data (CRITICAL)}

All results use synthetic test data for pipeline validation, NOT real GitHub commits. This is acceptable for a methodological paper demonstrating: (1) pipeline functionality, (2) permutation test methodology, (3) gate evaluation logic. Scientific validity for behavioral claims: ZERO.

Required next step: Collect 10K real GitHub commits. Real data collection (Phase 1A: 5 days mining, Phase 1B: 22 hours metric computation) is required before scientific conclusions can be drawn.

\textbf{L2: Commit Granularity}

When applied to real data, AST distance filter (< 20 nodes) will isolate small changes that may be too fine-grained to exhibit directional effects. This reduces confounds from multi-faceted large changes but tests whether even "clean" single-aspect commits show factorization.

Constructive future direction: Stratified analysis by commit size (20-50, 50-100 nodes) could reveal scale-dependent factorization.

\textbf{L3: Language and Domain Coverage}

Methodology targets Python/JavaScript commits. Systems languages (C++, Rust) or specialized domains may differ. Future work: Replicate across languages, measure cross-language consistency.

\textbf{L4: Metric Reliability}

Phase 0 metric validation (test-retest ICC ≥ 0.8, construct validity r ≥ 0.7) was not executed for synthetic data (not applicable). For real data: Execute Phase 0 validation before data collection, drop metrics with ICC < 0.8.

\textbf{L5: Label Quality}

Phase 1B purity validation (n=500 expert annotations, target > 70\% single-aspect) was not executed. Permutation test is robust to label noise by design. For real data: Expert annotation to measure purity. If 60-70\%, apply mixture modeling. If < 60\%, labels too noisy.

\subsubsection{Alternative Structural Hypotheses}

If global factorization fails in real data, structure may exist in alternative forms:

\textbf{H-A1: Hierarchical Quality Structure (DAG)}
\begin{itemize}
\item Claim: Aspects form dependency DAG: Correctness → Quality → Security → Efficiency
\item Test: Structural equation modeling with DAG priors, compare BIC/AIC vs PCA
\item Contribution: Would suggest causal ordering rather than simultaneous optimization
\end{itemize}

\textbf{H-A2: Contextual Factorization (Domain Clusters)}
\begin{itemize}
\item Claim: Aspect separation exists within domains (crypto repos for security) but not globally
\item Test: Mixture modeling with domain-specific covariance matrices
\item Contribution: Would justify domain-conditional routing rather than universal factorization
\end{itemize}

\textbf{H-A3: Commit Size Dependence (Scale-Dependent)}
\begin{itemize}
\item Claim: Large commits (50-100 nodes) exhibit factorization, minimal commits (< 20) don't
\item Test: Stratify by AST distance, test gap per stratum
\item Contribution: Would inform architectural decisions based on modification scale
\end{itemize}

\textbf{H-A4: Temporal Dynamics (Sequence Analysis)}
\begin{itemize}
\item Claim: Aspect structure emerges over commit sequences, not snapshots
\item Test: Time-series analysis of commit trajectories, directional momentum
\item Contribution: Would suggest sequential optimization rather than simultaneous
\end{itemize}

\subsubsection{Broader Impact}

\textbf{Positive impacts if finding holds (contingent on real data validation)}:
\begin{itemize}
\item Research efficiency: Prevents wasted effort on architectures lacking empirical justification
\item Practitioner guidance: Validates simpler methods (weighted scalarization) as principled defaults
\item Transparency: Establishes validation standard—test separability before architectural commitment
\item Negative result value: High-impact negative results redirect fields efficiently
\end{itemize}

\textbf{No ethical concerns}:
\begin{itemize}
\item Analysis of public GitHub data (no private code)
\item No personally identifiable information
\item Findings about code structure, not individuals
\item No potential for misuse
\end{itemize}
